\chapter{The Daily Workflow}\label{ch:workflow}

\section{Solving a problem, step by step}

\begin{figure}[H]
\centering
\begin{tikzpicture}[
  stp/.style={draw=cpteal, fill=cpteal!7, rounded corners=4pt, thick, text width=3.35cm, minimum height=2.15cm,
    align=center, font=\sffamily\small},
  num/.style={circle, fill=cpteal, text=white, font=\sffamily\bfseries\small, inner sep=1.5pt, minimum size=5.5mm},
  arr/.style={->, very thick, cpteal!70}]
  \foreach \px/\ti/\tx [count=\k] in {
      0/Fetch/{Competitive Companion button on the problem page},
      4.1/Read/{problem text, constraints, sample I/O},
      8.2/Code/{\texttt{solve()}; paste templates with \key{Ctrl+Alt+T}},
      12.3/Test/{CPH \emph{Run All} \key{Ctrl+Alt+B}}} {
    \node[stp] (s\k) at (\px, 0) {\textbf{\ti}\\[2pt]\scriptsize \tx};
    \node[num] at (s\k.north west) {\k};
  }
  \foreach \px/\ti/\tx [count=\k from 5] in {
      12.3/Explain/{\key{F5}: crash reason, time, memory per test},
      8.2/Hunt/{\key{Ctrl+Alt+X}: stress test vs a brute force},
      4.1/Fix/{add the failing input to CPH, fix, rerun},
      0/Submit/{the file exactly as it is}} {
    \node[stp] (s\k) at (\px, -3.1) {\textbf{\ti}\\[2pt]\scriptsize \tx};
    \node[num] at (s\k.north west) {\k};
  }
  \draw[arr] (s1) -- (s2); \draw[arr] (s2) -- (s3); \draw[arr] (s3) -- (s4);
  \draw[arr] (s4) -- node[right, font=\sffamily\scriptsize, text=cpgray] {fails} (s5);
  \draw[arr] (s5) -- (s6); \draw[arr] (s6) -- (s7);
  \draw[arr] (s7) -- node[above, font=\sffamily\scriptsize, text=cpgray] {passes} (s8);
  \draw[arr, cpgreen] (s4.north) to[out=90, in=90, looseness=0.35] node[above, font=\sffamily\scriptsize, text=cpgreen] {all tests pass: go straight to submit} (s1.north);
\end{tikzpicture}
\caption{Steps 5--7 are only needed when a test fails and the reason is not obvious.}
\end{figure}

\subsection{1. Fetch}
Click the Competitive Companion button in the browser. CPH receives the problem (name, contest, link, samples, limits
and whether it reads from a file) and creates \file{A\_Problem\_Name.cpp} from \file{template/template.cpp}.
The file opens with its header already filled in:

\begin{code}
/**
 *    author:  Kshitish Madbhavi
 *    problem: A. Graph
 *    source:  Codeforces - Baltic Olympiad in Informatics 2020, Day 2 (IOI, Unofficial Mirror Contest, Unrated)
 *    url:     https://codeforces.com/contest/1387/problem/A
 *    created: 2026-09-29 20:31:55
 */
\end{code}

For a USACO problem the same step also switches on file input and output; you do not touch anything
(Section~\ref{sec:usaco}).

\subsection{2. Read}
Look at the constraints before thinking about an algorithm: they tell you the complexity you can afford
(ALGORITHMS.pdf, Chapter~2 has the full table). Decide whether the input has several test cases. If it starts
with ``the first line contains $t$'', uncomment \cc{cin >> t;} in \cc{main} (put the cursor on that line and press
\key{Ctrl+/}).

\subsection{3. Code}
Write the solution in \cc{solve()}. Three ways to pull in library code, all of which paste at the cursor:

\begin{center}
\begin{tabularx}{\linewidth}{p{4.3cm}L}
\toprule
\thd{Method} & \thd{Example} \\
\midrule
Chord shortcut & \key{Ctrl+Alt+T} then \key{J} pastes Dijkstra; \key{Ctrl+Alt+Y} then \key{H} pastes HLD. \\
Prefix + \key{Tab} & type \texttt{cpdsu} and press \key{Tab}. Every prefix starts with \texttt{cp}, so typing \texttt{cp} lists all of them. \\
Picker & \key{Ctrl+Alt+T} then \key{Space}: a searchable list of every template. \\
\bottomrule
\end{tabularx}
\end{center}

Chapter~\ref{ch:library} explains the library and has keyboard maps of every chord.

\subsection{4. Test}
Press \key{Ctrl+Alt+B} or CPH's \emph{Run All} button. CPH compiles with the flags from Chapter~\ref{ch:compiler}
(about 0.4--0.9\,s thanks to the precompiled header) and runs every sample. \cc{dbg} output appears in each test's
\emph{Standard Error} box without affecting the verdict. Add your own tests with CPH's \emph{New Testcase} button;
edge cases (smallest $n$, largest $n$, all equal values, negative numbers) catch most wrong answers.

\subsection{5--7. When a test fails}
\begin{itemize}
  \item \textbf{Wrong answer and you can see why}: fix it and rerun.
  \item \textbf{Runtime error, or a time you do not understand}: press \key{F5}. The same CPH tests run in the terminal
        and each failure shows the input, expected output, your output, your stderr and a plain-language crash
        reason such as ``stack overflow'' (Section~\ref{sec:f5}).
  \item \textbf{The samples pass but the judge says wrong answer}: write a slow brute force in \file{stress/brute.cpp}
        and a random input generator in \file{stress/gen.cpp}, then press \key{Ctrl+Alt+X}. The stress tester finds a
        small input where your answer differs and copies it to the clipboard; paste it into CPH as a new test
        (Section~\ref{sec:stress}).
\end{itemize}

\subsection{8. Submit}
Submit the file unchanged: \cc{dbg} lines compile to nothing on the judge, and USACO file I/O is already on. CPH's
\emph{Submit} (\key{Ctrl+Alt+S}) sends Codeforces submissions as \emph{GNU G++23 14.2 (64 bit)}. Your code is still
strictly C++17: that is what CPH and IntelliSense compile it as, and every library template is checked with
\texttt{-std=c++17 -Wpedantic}. G++23 was kept because it is the newest 64-bit compiler Codeforces offers and
accepts C++17 code unchanged.

\section{Keyboard shortcuts}

\begin{center}
\begin{tabularx}{\linewidth}{>{\raggedright\arraybackslash}p{4.4cm}L}
\toprule
\thd{Keys} & \thd{Action} \\
\midrule
\key{Ctrl+Alt+B} & CPH: compile and run all test cases (CPH's own shortcut). \\
\key{Ctrl+Alt+S} & CPH: submit to Codeforces. \\
\key{Ctrl+Alt+T}, \emph{key} & paste a core template (Bronze to Gold level). \\
\key{Ctrl+Alt+Y}, \emph{key} & paste an advanced template (Platinum and beyond). \\
\key{Ctrl+Alt+U}, \emph{key} & paste a specialist template (polynomials, matchings, cuts, exotic structures). \\
\key{Ctrl+Alt+T}, \key{Space} & searchable template picker. \\
\key{F5} & run the CPH test cases in the terminal with diagnostics. \\
\key{Ctrl+F5} & run with input you type in the terminal. \\
\key{Ctrl+Alt+X} & stress test the current file. \\
\key{Ctrl+Shift+N} & new problem file from the template (asks for a name). \\
\key{Ctrl+Shift+U} & new USACO problem (asks for the file name, e.g.\ \texttt{moocast}). \\
\key{Ctrl+Shift+Alt+C} & archive solutions untouched for a day into \file{archive/YYYY-MM/}. \\
\key{Ctrl+Alt+N} & Code Runner: same as \key{Ctrl+F5}. \\
\bottomrule
\end{tabularx}
\end{center}

\section{Multi-test problems}

Competitive Companion does not tell CPH whether a problem has several test cases (it always reports
``single''), so this cannot be detected automatically. The template therefore reads one test by default:

\begin{code}
int t = 1;
// cin >> t;
while (t--) solve();
\end{code}

For a multi-test problem, uncomment the middle line. Remember that everything \cc{solve()} uses must be reset
between test cases: declare arrays inside \cc{solve()} (sized by that test's $n$) rather than globally.

\begin{pitfall}[Resetting global state]
The most common multi-test bug is a global array or counter that keeps values from the previous test.
Either declare it inside \cc{solve()} or clear exactly the part you used; never \cc{memset} a $10^6$-element
array $10^5$ times, which is $10^{11}$ operations.
\end{pitfall}

\section{USACO problems}\label{sec:usaco}

USACO reads from \texttt{name.in} and writes \texttt{name.out}. When Competitive Companion sends a USACO problem it
tells CPH the file names, and CPH writes them into the template (Section~\ref{sec:fileio} shows exactly how).
On the USACO judge your program opens the files; locally, CPH and F5 keep feeding it standard input, so the
samples still run in CPH. There is nothing to uncomment.

For a USACO problem typed in by hand, press \key{Ctrl+Shift+U} and enter the file name (\texttt{moocast}): the new file
behaves exactly like one created by CPH.

\section{Quickstart and archiving}\label{sec:quickstart}

\begin{itemize}
  \item \key{Ctrl+Shift+N} asks for a name, turns it into a safe file name (\texttt{B. Two Sum} becomes
        \file{B\_Two\_Sum.cpp}), fills the template exactly as CPH would and opens the file. Use it for problems
        without Competitive Companion support. If the file exists it is simply opened.
  \item \key{Ctrl+Shift+U} does the same for USACO and fills in \texttt{name.in} / \texttt{name.out}.
  \item \key{Ctrl+Shift+Alt+C} moves every solution in the folder root that has not been modified for a day into
        \file{archive/YYYY-MM/} (by the month it was last edited), deletes stray \texttt{.exe} / \texttt{.bin} files,
        and moves each problem's CPH test cases along with it. CPH names a problem's test file after the MD5 hash of
        the source path, so the test file is renamed for the new path and CPH still finds your tests if you open the
        archived file. Nothing is ever deleted except build output.
\end{itemize}
